\documentclass[10pt,a4paper]{amsart}
\usepackage[margin=19mm]{geometry}
\usepackage{amsmath,amssymb,mathtools}
\usepackage[scr=rsfs,cal=euler]{mathalfa}
\usepackage{xfrac}
\usepackage[unicode,hidelinks,bookmarks=false]{hyperref}
\newcommand{\ie}{i.e.\ }
\newcommand{\cf}{cf.\ }
\newcommand{\resp}{respectively\ }
\newcommand{\Subsection}{\subsection}
\providecommand{\nobreakdash}{\nobreak\hbox{-}}
\setlength{\emergencystretch}{2em}
\multlinegap0pt
\usepackage[all]{xy}
\usepackage{amssymb}
\usepackage{xcolor}
\newtheorem{thm}{Theorem}
\newtheorem{prop}{Proposition}
\title{Arithmetic BF theory and the Cassels-Tate pairing}
\author{Jeehoon Park, Junyeong Park}
\date{Source: arXiv:2602.19621v2 (CC BY 4.0)}
\begin{document}
\maketitle

\noindent\emph{Source and licence.} Arithmetic BF theory and the Cassels-Tate pairing. Version arXiv:2602.19621v2 (CC BY 4.0), distributed by arXiv under the Creative Commons Attribution 4.0 International licence, http://creativecommons.org/licenses/by/4.0/, as recorded in the arXiv licence metadata for this exact version and shown on the abstract page https://arxiv.org/abs/2602.19621v2. Full licence notice: \texttt{licenses/arxiv-2602.19621-CC-BY-4.0.txt}.\par\smallskip

\noindent\textit{Editorial scope.} Counted unit: the article's prop ctpbf (source0.tex lines 1693--1704). The article's complete original proof of it is delivered with it (source0.tex lines 1705--1741, one \texttt{proof} environment of 37 lines). No mathematics is written, repaired or re-derived: the statement and the proof are the article's own lines, in the article's own notation, and no retained line is substituted or re-flowed.  Retained uncounted as the source-local context the proof needs: lines 782--784; lines 801--803; lines 891--898; lines 1057--1059; lines 1152--1163.  Nothing was omitted from any retained span, so the package carries no omission record.  No retained line contains a citation, so the wrapper carries no bibliography and no bibliography entry.  No retained line contains a figure environment, an image-inclusion command or drawing code, so no image is delivered and the package declares no asset dependency.  Editorial formatting only, in addition: an AMS \texttt{amsart} preamble that restates the article's own declarations and macros used by the retained lines, namely the environments \texttt{thm}, \texttt{prop} on separate counters, together with the standard packages those lines need.  The retained environments print in this document as Theorem 1, Proposition 1 in order of appearance, whereas the article prints the same items with the numbering of its own full document.  The complete native source of the article as distributed in the arXiv e-print for this exact version is bundled unchanged as \texttt{sources/195128/source0.tex}.
\begin{align}\label{local-HGnr-vanishing}
    H^q(G_{K,x}^\mathrm{nr},\mathcal{O}_x^{\mathrm{nr},\times})=0\quad\textrm{for $q\geq2$.}
\end{align}

\begin{align}\label{BFnrx-lift}
    d\varphi_x^\mathrm{nr}=\alpha_{x,1}^\mathrm{nr}\cup d(\sigma\circ\alpha_{x,2}^\mathrm{nr})\quad\textrm{for some}\quad\varphi_x^\mathrm{nr}\in\frac{C^2(G_{K,x}^\mathrm{nr},\mathcal{O}_x^{\mathrm{nr},\times})}{B^2(G_{K,x}^\mathrm{nr},\mathcal{O}_x^{\mathrm{nr},\times})}\rlap{\ .}
\end{align}

\begin{align}\label{FXS-defn}
\begin{aligned}
    \xymatrix{
    \mathscr{F}(X_S) \ar[d] \ar[r] & H^1(G_K,M_1^\vee)\times H^1(G_K,M_2) \ar[d] \\
    \mathrm{BC}(\partial X_S) \ar[r] & \displaystyle\prod_{x\in X^\mathrm{cl}}\mathscr{F}_x\rlap{\ .}
    }
\end{aligned}
\end{align}

\begin{align}\label{FUS-BF-functional}
    \xymatrix{\mathrm{BF}_{X_S}:\mathscr{F}(X_S) \ar[r] & \mathscr{L}_S & \rho \ar@{|->}[r] & [(\partial_x\varphi\bmod\Theta_x)_{x\in S}]}
\end{align}

\begin{thm}[Decomposition formula]\label{Decomp-formula} Let $S\subseteq T\subseteq X^\mathrm{cl}$ be finite subsets.
\begin{quote}
    (1) If $S\neq\emptyset$, then the following equality holds in $\Gamma(\mathscr{F}(X_S),\partial_T^\ast\mathscr{L}_T)$:
    \begin{align*}
        \mathrm{BF}_{X_T}|_{\mathscr{F}(X_S)}=\mathrm{BF}_{X_S}\oplus\left(\mathrm{BF}_{T\setminus S}^\mathrm{nr}\circ\partial^\mathrm{nr}_{T\setminus S}\right)\rlap{\ .}
    \end{align*}
    (2) If $S=\emptyset$, $X\setminus Y\subseteq T$, and the order of $\mathcal{M}$ is invertible on $Y_T$, then the following equality holds in $\Gamma(\mathscr{F}(X),\partial_T^\ast\mathscr{L}_T)$:
    \begin{align*}
        \mathrm{BF}_{X_T}|_{\mathscr{F}(X)}=\mathrm{BF}_X\oplus\left(\mathrm{BF}_T^\mathrm{nr}\circ\partial_T^\mathrm{nr}\right)\rlap{\ .}
    \end{align*}
\end{quote}
\end{thm}

\begin{prop} \label{ctpbf}
Assume that
\begin{eqnarray}
    \iota^{-1}(\mathscr{W}_y)^{\perp}\times \pi(\mathscr{W}_y)= \mathscr{F}_y^\mathrm{nr}
    %:=H^1(G_{K,x}^\mathrm{nr},(M_1^\vee)^{I_{K,x}})\times H^1(G_{K,x}^\mathrm{nr},\pi(M^{I_{K,x}})),
    \quad y \in Y^{\mathrm{cl}}.
\end{eqnarray}
Then $\mathrm{Sel} (M_2,\mathscr{W}_2) \times \mathrm{Sel} (M_1^\vee, \mathscr{W}_1^\perp)$ is identified with $\mathscr{F}(X)$ and we have
    \begin{eqnarray*}
        \mathrm{CTP}_E(\rho_1, \rho_2)= \mathrm{BF}_X(\rho_1, \rho_2).
    \end{eqnarray*}
\end{prop}

\begin{proof}
Given $\rho=(\rho_1,\rho_2)\in\mathscr{F}(X)$, choose a representative of its image under \eqref{FXS-defn}:
\begin{align*}
    \alpha=(\alpha_1,\alpha_2)\in Z^1(G_K,M_1^\vee)\times Z^1(G_K,M_2)
\end{align*}


By \eqref{BFnrx-lift}, for each $x \in Y^{\mathrm{cl}}$,
\begin{align*}
    d\varphi_x^\mathrm{nr}=\partial_x\alpha_{1}\cup d(\sigma\circ\partial_x\alpha_{2})\quad\textrm{for some}\quad\varphi_x^\mathrm{nr}\in\frac{C^2(G_{K,x}^\mathrm{nr},\mathcal{O}_x^{\mathrm{nr},\times})}{B^2(G_{K,x}^\mathrm{nr},\mathcal{O}_x^{\mathrm{nr},\times})}\rlap{\ .}
\end{align*}
For each $x \in X^{\mathrm{cl}}$, we have
\begin{eqnarray*}
    d\left( \partial_x\alpha_{1} \cup( \alpha_{x,2,M}-\sigma \circ \partial_x\alpha_{2} )\right) = \partial_x\alpha_{1}\cup d(\sigma\circ \partial_x\alpha_{2}).
\end{eqnarray*}
Thus, \eqref{local-HGnr-vanishing} implies that
\begin{eqnarray}\label{sectionunr}
    \partial_x\alpha_{1} \cup( \alpha_{x,2,M}-\sigma \circ \partial\alpha_{2}) - \varphi_x^{\mathrm{nr}} \in B^2(G_{K,x}, \mathscr{O}_x^{\mathrm{nr}, \times}) 
    %= Z^2(G_{K,x}, \mathscr{O}_x^{\mathrm{nr}, \times})
\end{eqnarray}
is a coboundary.
Note that $d\epsilon = \alpha_1\cup d f  $ and we choose $\sigma \in \mathscr{R}$ such that $f=\sigma \circ \alpha_2$ and choose $\varphi \in C^2(G_K, K^{\mathrm{sep},\times})/B^2(G_K, K^{\mathrm{sep},\times})$ (as in \eqref{FUS-BF-functional}) such that 
\begin{eqnarray}\label{globalchoice}
   d \varphi = \alpha_1 \cup d(\sigma \circ \alpha_2)=d\epsilon. 
\end{eqnarray}


Then we have the following equality in $\mathbb{Q}/\mathbb{Z}$:
\begin{eqnarray*}
    \mathrm{CTP}_E(\rho_1, \rho_2) &:=& \sum_{v \in X^{\mathrm{cl}}} \mathrm{inv}_v (\gamma_v)\\
   &=& \sum_{v \in X^{\mathrm{cl}}} \mathrm{inv}_v \left(\partial_x\epsilon- \partial_x\alpha_{1} \cup( \alpha_{x,2,M}-\sigma \circ \partial_x\alpha_{2} ) \right)\\
   &\stackrel{\eqref{globalchoice}}{=}& \sum_{v \in X^{\mathrm{cl}}} \mathrm{inv}_v \left(\partial_x\varphi- \partial_x\alpha_{1}\cup ( \alpha_{x,2,M} -\sigma \circ \partial_x\alpha_{2} )\right)\\
   &\stackrel{\eqref{sectionunr}}{=}&\sum_{x\in T}\mathrm{inv}_x(\partial_x\varphi-\varphi_x^{\mathrm{nr}})\quad \text{ for some finite set $T$ of places}\\
   &=&\mathrm{BF}_{X_T}(\rho)-\mathrm{BF}_T^\mathrm{nr}(\partial^\mathrm{nr}_T\rho)\rlap{\ }
   \stackrel{\mathrm{Theorem}\ \ref{Decomp-formula}, (2)}{=} \mathrm{BF}_X(\rho_1, \rho_2).
\end{eqnarray*}
\end{proof}

\end{document}
